% Part: turing-machines
% Chapter: undecidability
% Section: enumerating-tms

\documentclass[../../../include/open-logic-section]{subfiles}

\begin{document}

% SEGMENT OLP-0266-S01

\olfileid{tur}{und}{enu}
\olsection{டியூரிங் பொறிகளைப் பட்டியலிடுதல்}

\begin{explain}
எல்லா டியூரிங் பொறிகளின் கணமும் !!{enumerable} என்று காட்டலாம்.
ஒவ்வொரு டியூரிங் பொறியையும் முடிவுறுமாறு விவரிக்கலாம் என்ற உண்மையிலிருந்து
இது பின்வருகிறது. நிலைகளின் கணமும் நாடாவின் குறித்தொகுதியும் முடிவுறு
கணங்கள். நிலைமாற்றுச் சார்பு என்பது $Q \times \Sigma$ இலிருந்து
$Q \times \Sigma \times \{\TMleft, \TMright, \TMstay\}$ குச் செல்லும்
ஒரு பகுதிச் சார்பு; எனவே அது வரையறுக்கப்பட்டிருக்கும் முடிவுறு
எண்ணிக்கையிலான மாறி ஜோடிகளுக்கான அதன் மதிப்புகளைப் பட்டியலிட்டும்
அதைக் குறிப்பிடலாம்.

இவ்வளவு வரை இது உண்மை; ஆனால் ஒரு நுணுக்கமான வேறுபாடு உள்ளது.
நிலைகளின் கணமும் நாடாவின் குறியெழுத்துத் தொகுதியும் எந்த
!!{element}s ஐக் கொண்டிருக்கலாம் என்பதில் டியூரிங் பொறிகளின் வரையறை
எந்தக் கட்டுப்பாட்டையும் விதிக்கவில்லை. ஆகவே, எடுத்துக்காட்டாக,
ஒவ்வொரு மெய்யெண்ணுக்கும் அந்த எண்ணை ஒரு நிலையாகப் பயன்படுத்தும்
டியூரிங் பொறி ஒன்று தொழில்நுட்பமுறையில் உள்ளது. எனினும், எந்தப்
பொருள்கள் நிலைகளாகவும் குறித்தொகுதியாகவும் செயல்படுகின்றன என்பதிலிருந்து
டியூரிங் பொறியின் \emph{நடத்தை} சார்பற்றது. \olref{fig:variants}
இலுள்ள இரு டியூரிங் பொறிகளைக் கருதுக.
\begin{figure}
\begin{center}
\begin{tikzpicture}[->,>=stealth',shorten >=1pt,auto,node distance=2.8cm,
                    semithick]
  \tikzstyle{every state}=[fill=none,draw=black,text=black]

  \node[initial,state]         (A)              {$q_0$};
  \node[state]         (B) [right of=A] {$q_1$};

  \path (A) edge [bend left] node {\TMtrans{\TMstroke}{\TMstroke}{\TMright}} (B)
        (B) edge [loop above] node {\TMtrans{\TMblank}{\TMblank}{\TMright}} (B)
            edge [bend left] node {\TMtrans{\TMstroke}{\TMstroke}{\TMright}} (A);
\end{tikzpicture}\\
\begin{tikzpicture}[->,>=stealth',shorten >=1pt,auto,node distance=2.8cm,
  semithick]
\tikzstyle{every state}=[fill=none,draw=black,text=black]

\node[initial,state]         (A)              {$s$};
\node[state]         (B) [right of=A] {$h$};

\path (A) edge [bend left] node {\TMtrans{A}{A}{\TMright}} (B)
(B) edge [loop above] node {\TMtrans{\TMblank}{\TMblank}{\TMright}} (B)
edge [bend left] node {\TMtrans{A}{A}{\TMright}} (A);
\end{tikzpicture}
\end{center}
\caption{\emph{இரட்டை எண்} பொறியின் மாறுவகைகள்}
\ollabel{fig:variants}
\end{figure}
இந்த இரு வரைபடங்களும் இரு பொறிகளுக்கு ஒத்துள்ளன: நாடாவின்
குறியெழுத்துத் தொகுதி $\Sigma = \{\TMendtape,\TMblank,\TMstroke\}$
மற்றும் நிலைகளின் கணம் $\{q_0,q_1\}$ ஆகியவற்றைக் கொண்ட~$M$;
குறியெழுத்துத் தொகுதி $\Sigma' =
\{\TMendtape,\TMblank,A\}$ மற்றும் நிலைகள் $\{s,h\}$ ஆகியவற்றைக்
கொண்ட~$M'$. ஆனால் மற்ற வகைகளில் அவற்றின் கட்டளைகள் ஒரே மாதிரியானவை:
$n$ இரட்டை எண் எனில் மட்டுமே,~$M$ ஆனது $n$ எண்ணிக்கையிலான
$\TMstroke$-களின் ஒரு தொடரில் நிற்கும்; $n$ இரட்டை எண் எனில் மட்டுமே,
~$M'$ ஆனது $n$ எண்ணிக்கையிலான $A$-களின் ஒரு தொடரில் நிற்கும். நாம்
செய்ததெல்லாம் $\TMstroke$ ஐ~$A$ என்றும், $q_0$ ஐ~$s$ என்றும், $q_1$
ஐ~$h$ என்றும் மறுபெயரிட்டதே. இவ்வெடுத்துக்காட்டைப் பொதுமைப்படுத்தலாம்:
குறிகளையும் நிலைகளையும் இவ்வாறு மறுபெயரிடுவதன் மூலம் ஒன்றிலிருந்து
மற்றொன்று கிடைக்கும் வரை டியூரிங் பொறிகளை ஒரே பொறியாகக் கருதலாம்.
உண்மையில், ஒரு டியூரிங் பொறியின் குறிகளையும் நிலைகளையும் நேர்ம
முழுக்களாகவே எளிதில் கருதலாம்: $\sigma_0$ குப் பதிலாக~$1$ ஐயும்,
$\sigma_1$ குப் பதிலாக~$2$ ஐயும், இவ்வாறாகத் தொடர்ந்தும் கருதுக;
$\TMendtape$ என்பது~$1$, $\TMblank$ என்பது~$2$, இவ்வாறாகத் தொடரும்.
இவ்வழியில், \emph{இரட்டை எண்} பொறி \olref{fig:standard-even} இல்
படம்பிடிக்கப்பட்ட பொறியாகிறது.
\begin{figure}
\[\begin{tikzpicture}[->,>=stealth',shorten >=1pt,auto,node distance=2.8cm,
  semithick]
\tikzstyle{every state}=[fill=none,draw=black,text=black]

\node[initial,state]         (A)              {$1$};
\node[state]         (B) [right of=A] {$2$};

\path (A) edge [bend left] node {\TMtrans{3}{3}{\TMright}} (B)
(B) edge [loop above] node {\TMtrans{2}{2}{\TMright}} (B)
edge [bend left] node {\TMtrans{3}{3}{\TMright}} (A);
\end{tikzpicture}
\]
\caption{ஒரு தரநிலை \emph{இரட்டை எண்} பொறி}
\ollabel{fig:standard-even}
\end{figure}
நேர்ம முழுக்களை நிலைகளாகவும் குறிகளாகவும் கொண்ட ஒரு டியூரிங் பொறியை
\emph{தரநிலைப்} பொறி என்று அழைத்து, இனிமேல் தரநிலைப் பொறிகளை மட்டும்
கருதலாம்.\footnote{``தரநிலைப் பொறி'' என்ற சொல்வழக்கு தரநிலையானது அல்ல.}

டியூரிங் பொறிகளின் கணம் !!{enumerable} என்று காட்ட விரும்பினோம்.
மேலுள்ள கருத்துகளை மனத்தில் கொண்டு, தரநிலை டியூரிங் பொறிகளின் கணம்
!!{enumerable} என்று காட்டினால் போதும். ஒரு தரநிலை டியூரிங் பொறி
$M = \tuple{Q, \Sigma, q_0, \delta}$ நமக்குக் கொடுக்கப்பட்டுள்ளது
என்று கொள்க. நேர்ம முழுக்களின் ஒரு முடிவுறு சரத்தைப் பயன்படுத்தி அதை
எவ்வாறு விவரிக்கலாம்? முதலில் நிலைகளின் எண்ணிக்கை, நிலைகள், குறிகளின்
எண்ணிக்கை, குறிகள், தொடக்க நிலை ஆகியவற்றைப் பட்டியலிடுவோம். ($M$ ஒரு
தரநிலைப் பொறி என்பதால் இவை அனைத்தும் நேர்ம முழுக்கள் என்பதை நினைவில்
கொள்க.)~$\delta$ பற்றி என்ன? $Q$,~$\Sigma$ ஆகியவை முடிவுறுவதால்,
சாத்திய மாறிகளின் கணம், அதாவது ஜோடிகள் $\tuple{q,\sigma}$, முடிவுறும்.
எனவே~$\delta$ இலுள்ள தகவல், $\delta(q,\sigma) =
\tuple{q', \sigma', D}$ ஆக உள்ள எல்லா $5$-உறுப்படிகள்
$\tuple{q, \sigma, q', \sigma', d}$ இன் முடிவுறு பட்டியல் மட்டுமே;
இங்கு $d$ என்பது திசை~$D$ ஐக் குறியீடாக்கும் ஓர் எண் (எடுத்துக்காட்டாக,
~$\TMleft$ கு $1$,~$\TMright$ கு $2$,~$\TMstay$ கு $3$).

இவ்வழியில், ஒவ்வொரு தரநிலை டியூரிங் பொறியையும் நேர்ம முழுக்களின் ஒரு
முடிவுறு பட்டியலால், அதாவது $s_M \in (\PosInt)^*$ என்ற ஒரு தொடராக
விவரிக்கலாம். எடுத்துக்காட்டாக, தரநிலை \emph{இரட்டை எண்} பொறி
பின்வரும் தொடரால் குறியீடாக்கப்படுகிறது:
\[
2, \underbrace{1, 2}_Q, 3, \overbrace{1, 2, 3}^\Sigma, 1, \underbrace{1, 3, 2, 3, 2}_{\delta(1,3) = \tuple{2,3,R}}, 
\overbrace{2, 2, 2, 2, 2}^{\delta(2,2) = \tuple{2,2,R}},
\underbrace{2, 3, 1, 3, 2}_{\delta(2,3) = \tuple{1,3,R}}.
\]
\end{explain}

\begin{thm}
$\Nat$ இலிருந்து~$\Nat$ குச் செல்லும், டியூரிங்-கணிக்கத்தக்கவையல்லாத
சார்புகள் உள்ளன.
\end{thm}

\begin{proof}
நேர்ம முழுக்களின் முடிவுறு தொடர்களின் கணம்~$(\PosInt)^*$
!!{enumerable} என்பதை அறிவோம்
% READER FALLBACK TA-READER-REF-011: retain the live SFR exercise link when present.
(\oliflabeldef{sfr:siz:zigzag:prob:posint-star}
  {\cref{sfr:siz:zigzag:prob:posint-star}}
  {நேர்ம முழுக்களின் முடிவுறு தொடர்கள் பட்டியலிடத்தக்கவை என நிறுவும் பயிற்சி}).
தரநிலை டியூரிங் பொறிகளின்
விவரிப்புகளின் கணம்~$(\PosInt)^*$ இன் ஒரு உட்கணமாக இருப்பதால், அதுவும்
பட்டியலிடத்தக்கது என்பது இதிலிருந்து கிடைக்கிறது. $\Nat$ இலிருந்து
~$\Nat$ குச் செல்லும் ஒவ்வொரு டியூரிங்-கணிக்கத்தக்க சார்பையும் ஏதோ
ஒரு டியூரிங் பொறி (உண்மையில், பல பொறிகள்) கணிக்கிறது. அதன் நிலைகளையும்
குறிகளையும் நேர்ம முழுக்களாக மறுபெயரிடுவதன் மூலம் (குறிப்பாக,
$\TMendtape$ ஐ~$1$ என்றும், $\TMblank$ ஐ~$2$ என்றும், $\TMstroke$
ஐ~$3$ என்றும்) ஒவ்வொரு டியூரிங்-கணிக்கத்தக்க சார்பையும் ஒரு தரநிலை
டியூரிங் பொறி கணிக்கிறது என்று காணலாம். $\Nat$ இலிருந்து~$\Nat$ குச்
செல்லும் எல்லா டியூரிங்-கணிக்கத்தக்க சார்புகளின் கணமும்
பட்டியலிடத்தக்கது என்பது இதன் பொருள்.

மறுபுறம், $\Nat$ இலிருந்து~$\Nat$ குச் செல்லும் எல்லாச் சார்புகளின்
கணம் !!{enumerable} அல்ல
% READER FALLBACK TA-READER-REF-011: retain the live SFR exercise link when present.
(\oliflabeldef{sfr:siz:red:prob:nat-nat}
  {\cref{sfr:siz:red:prob:nat-nat}}
  {இயல் எண்களிலிருந்து இயல் எண்களுக்குச் செல்லும் எல்லாச் சார்புகளின் கணம் பட்டியலிடத்தக்கதல்ல என நிறுவும் பயிற்சி}).
எல்லாச்
சார்புகளும் ஏதோ ஒரு டியூரிங் பொறியால் கணிக்கத்தக்கவையாக இருந்தால்,
அவற்றைக் கணிக்கும் டியூரிங் பொறிகளின் எல்லா விவரிப்புகளையும்
பட்டியலிடுவதன் மூலம் எல்லாச் சார்புகளின் கணத்தையும் பட்டியலிடலாம்.
ஆகவே டியூரிங்-கணிக்கத்தக்கவையல்லாத சில சார்புகள் உள்ளன.
\end{proof}

\begin{prob}
  நிலைகளையும் குறியெழுத்துத் தொகுதியின் குறிகளையும் வெளிப்படையாகப்
  பட்டியலிடத் தேவையில்லாமல் டியூரிங் பொறிகளை விவரிக்கும் ஒரு வழியை
  உங்களால் எண்ணிப்பார்க்க முடியுமா? ``தரநிலைப்'' பொறி பற்றிய உங்கள்
  சொந்தக் கருத்தை வரையறுக்கலாம்; ஆனால் உங்கள் புதிய பொருளில் ஒவ்வொரு
  டியூரிங் பொறியையும் ஒரு ``தரநிலைப்'' பொறியால் ஏன் கணிக்க முடியும்
  என்பது பற்றி ஏதாவது கூறுக.
\end{prob}

\end{document}
